\chapter{Logica proposizionale classica: sintassi, semantica, deduzione naturale}

\textit{Lucidi: \texttt{lucidi/3.Propositional\_Logic} (con le annotazioni a
mano del prof). Approfondimento: Van Dalen, \emph{Logic and Structure}, §1.1,
§1.2, §1.4 (Natural Deduction), §1.5 (Completeness); Chang--Lee, cap.~2.}

\nt{
  Sommario della lezione: (1) sintassi, (2) semantica, (3) forme normali,
  (4) deduzione e refutazione, (5) deduzione naturale. Le prime slide
  (proposizioni, connettivi e i loro problemi, alfabeto, definizione induttiva
  di $\PROP$, principio di induzione) ripetono quanto già visto nel capitolo
  precedente: qui riporto solo ciò che è nuovo, comprese le note scritte a mano.
  Il prof scrive a mano \textbf{``proposition = you can ask if it is true''}.
}

% =====================================================================
\section{Sintassi e semantica di un sistema formale}
% =====================================================================

Per \emph{qualsiasi} sistema formale (linguaggio di programmazione, logica,
\dots) si distinguono:

\dfn{Sintassi e semantica}{
  \begin{itemize}
    \item \textbf{Sintassi}: le regole che dicono come si costruiscono le frasi
          \textbf{ben formate}. Si può dare con definizioni informali, con una
          \textbf{BNF}, con \textbf{definizioni induttive}, \dots
    \item \textbf{Semantica}: le regole che dicono qual è il
          \textbf{significato} delle frasi ben formate. Esistono semantiche
          operazionali, denotazionali, logiche, \dots
  \end{itemize}
}

Per le logiche:
\begin{itemize}
  \item \textbf{sintassi}: quali simboli rappresentano proposizioni, variabili,
        funzioni, predicati, e quali i connettivi, in modo da ottenere le
        \textbf{formule ben formate} (wff);
  \item \textbf{semantica}, due grandi ``regni'':
    \begin{itemize}
      \item \textbf{model theory}: che cosa è \textbf{vero};
      \item \textbf{proof theory}: che cosa è \textbf{dimostrabile};
      \item \emph{si spera che le due coincidano} (è il teorema di completezza,
            in fondo al capitolo).
    \end{itemize}
\end{itemize}

% =====================================================================
\section{Sintassi della logica proposizionale}
% =====================================================================

L'alfabeto è quello della Definizione 1.1.1 di Van Dalen (capitolo
precedente). Annotazioni del prof sulla slide:
\begin{itemize}
  \item come simboli proposizionali usiamo anche $A, B, C, \dots$;
  \item $\to$ è l'\textbf{implicazione materiale};
  \item oltre a $\bot$ (falso) usiamo $\top$ (\textbf{true}).
\end{itemize}

La stessa sintassi viene data in tre modi equivalenti.

\dfn{Sintassi I: definizione ricorsiva delle formule ben formate}{
  \begin{enumerate}[start=0]
    \item (aggiunta a mano) se $S$ è una formula, $(S)$ è una formula;
    \item un atomo è una formula;
    \item se $S$ è una formula, $\neg S$ è una formula (\textbf{negazione});
    \item se $S_1$ e $S_2$ sono formule, $S_1 \wedge S_2$ è una formula
          (\textbf{congiunzione});
    \item se $S_1$ e $S_2$ sono formule, $S_1 \vee S_2$ è una formula
          (\textbf{disgiunzione});
    \item tutte le formule ben formate si ottengono applicando le regole sopra.
  \end{enumerate}
  Abbreviazione: $S_1 \to S_2$ si può scrivere come $\neg S_1 \vee S_2$.
}

\dfn{Sintassi II: Backus--Naur Form (BNF)}{
  \[
  \begin{array}{rcl}
    \langle\text{formula}\rangle & ::= & \text{Atomic Proposition} \\
      & \mid & \neg \langle\text{formula}\rangle \\
      & \mid & \langle\text{formula}\rangle \wedge \langle\text{formula}\rangle \\
      & \mid & \langle\text{formula}\rangle \vee \langle\text{formula}\rangle \\
      & \mid & \langle\text{formula}\rangle \to \langle\text{formula}\rangle \\
      & \mid & \langle\text{formula}\rangle \leftrightarrow \langle\text{formula}\rangle \\
      & \mid & (\langle\text{formula}\rangle)
  \end{array}
  \]
  Si legge: ``$::=$'' = \emph{is}, ``$\mid$'' = \emph{or}.
}

\ex{BNF dei numeri naturali (lavagna)}{
  \[
    \langle\text{natural}\rangle ::= 0 \;\mid\; \langle\text{natural}\rangle\, S
  \]
  dove $S$ è il successore. Il numero $3$ è $(((0)S)S)S$: tre applicazioni
  della seconda produzione. È lo stesso schema ``caso base + passo'' della
  definizione induttiva di $\PROP$.
}

\textbf{Sintassi III: definizione induttiva} (Van Dalen 1.1.2): $\PROP$ è il più
piccolo insieme $X$ chiuso rispetto ad atomi, connettivi binari e negazione
(vedi capitolo precedente). È la forma \textbf{utile per dimostrare per
induzione} proprietà delle formule (Teorema 1.1.3, dove $\square$ indica un
connettivo binario qualsiasi).

\subsection{Sintassi concreta e sintassi astratta}

\dfn{Sintassi concreta e astratta}{
  \begin{itemize}
    \item La \textbf{sintassi concreta} è la stringa che scriviamo, ad esempio
          $A \wedge B \vee C$.
    \item La \textbf{sintassi astratta} è l'\textbf{albero} della formula: i
          nodi interni sono connettivi, le foglie atomi.
  \end{itemize}
  Una stringa senza parentesi può corrispondere a più alberi: servono
  \textbf{precedenze} e \textbf{associatività} per sceglierne uno solo.
}

\ex{$A \wedge B \vee C$}{
  \begin{center}
  \begin{tikzpicture}[level distance=8mm, sibling distance=12mm, baseline=(r.base)]
    \node (r) {$\vee$}
      child {node {$\wedge$} child {node {$A$}} child {node {$B$}}}
      child {node {$C$}};
  \end{tikzpicture}
  \qquad\qquad
  \begin{tikzpicture}[level distance=8mm, sibling distance=12mm, baseline=(r.base)]
    \node (r) {$\wedge$}
      child {node {$A$}}
      child {node {$\vee$} child {node {$B$}} child {node {$C$}}};
  \end{tikzpicture}
  \end{center}
  L'albero giusto è il \textbf{primo}, cioè $(A \wedge B) \vee C$, perché
  $\wedge$ lega più di $\vee$. Il secondo, $A \wedge (B \vee C)$, sulla lavagna
  è \textbf{barrato}.
}

\dfn{Precedenza e associatività (lavagna)}{
  \begin{itemize}
    \item \textbf{Precedenza}:
      \[ \neg \;>\; \wedge \;>\; \vee \;>\; \to \;=\; \leftrightarrow \]
    \item \textbf{Associatività}: $\wedge, \vee, \to, \leftrightarrow$ sono
          \textbf{associativi a destra}.
  \end{itemize}
}

\nt{
  Quindi $A \to B \to C$ significa $A \to (B \to C)$ (lo useremo negli esercizi
  di deduzione naturale) e $\neg A \wedge B$ significa $(\neg A) \wedge B$. Per
  $\wedge$ e $\vee$ l'associatività è irrilevante semanticamente (sono
  associativi), per $\to$ no: $(A \to B) \to C \not\equiv A \to (B \to C)$.
  Una slide successiva elenca il ``rango decrescente''
  $\leftrightarrow,\ \to,\ \vee,\ \wedge,\ \neg$ (dal più debole al più forte),
  coerente con la lavagna salvo che lì $\to$ e $\leftrightarrow$ sono alla
  pari.
}

\ex{Formule ben formate e non (lavagna)}{
  \begin{itemize}
    \item $\big((\neg\bot) \vee (A \wedge C)\big) \to \big(B \leftrightarrow
          (\neg\top)\big)$ è ben formata.
    \item $A \wedge \vee\, (\neg B)$ \textbf{non} è ben formata (barrata): tra
          $\wedge$ e $\vee$ manca una formula.
  \end{itemize}
}

% =====================================================================
\section{Semantica I (classica): modelli, cioè tavole di verità}
% =====================================================================

\nt{
  Il prof aggiunge a mano \textbf{``CLASSICAL''} al titolo: quella che segue è
  la semantica della logica \emph{classica} (a due valori). A lavagna scrive la
  conseguenza logica come $\Gamma \Vdash F$, con $\Gamma$ insieme
  \textbf{finito} di formule, e annota che la restrizione al caso finito è
  \textbf{``not standard!!!''}: normalmente $\Gamma$ può essere infinito (come
  in Van Dalen).
}

Per definire il significato di una formula serve il significato dei simboli
proposizionali: ``It is raining'' è vera o falsa? Dipende dal ``mondo''.
Serve la nozione di \textbf{interpretazione}.

\dfn{Interpretazione}{
  Data una formula proposizionale $G$, sia $\{A_1, \dots, A_n\}$ l'insieme
  degli atomi (a mano: \textbf{variabili}) che compaiono in $G$.
  Un'\textbf{interpretazione} $I$ di $G$ è un assegnamento di valori di verità
  ad $A_1, \dots, A_n$.
}

\ex{}{
  $G \triangleq (P \vee Q) \wedge \neg(P \wedge Q)$. Insieme delle variabili:
  $\{P, Q\}$. \emph{Una} interpretazione per $G$: $I = \{P = \mathbf{T},\
  Q = \mathbf{F}\}$.
}

\begin{itemize}
  \item Ogni atomo $A_i$ può valere \textbf{True} o \textbf{False}, mai
        entrambi \hfill (a mano: è questo che rende la logica \textbf{classica}).
  \item Data un'interpretazione $I$, una formula $G$ è \textbf{vera in $I$} sse
        $G$ viene valutata a True in $I$.
  \item Una formula con $n$ atomi distinti ha $2^n$ interpretazioni distinte.
  \item \textbf{Convenzione}: $\{P, \neg Q, \neg R, S\}$ rappresenta
        l'interpretazione $I : \{P = T,\ Q = F,\ R = F,\ S = T\}$.
\end{itemize}

\dfn{Modello}{
  Se $G$ è vera sotto l'interpretazione $I$ diciamo che $I$ è un
  \textbf{modello} di $G$ e scriviamo
  \[ I \models G. \]
  Si usa lo \textbf{stesso simbolo} $\models$ della conseguenza logica.
}

\nt{
  Un'interpretazione è la ``valutazione ristretta agli atomi'' del capitolo
  precedente; per il lemma di dipendenza dai soli atomi presenti, tanto basta.
  I due usi di $\models$ si distinguono da cosa c'è a sinistra:
  un'interpretazione ($I \models G$: $G$ è vera in $I$) oppure formule
  ($F_1, \dots, F_n \models G$: $G$ è vera in tutti i modelli delle $F_i$).
}

\subsection{Semantica dei connettivi}

Il significato dei connettivi è definito dalle \textbf{tavole di verità}: dati
i valori degli argomenti danno il valore della formula composta. In parole
(a mano: \textbf{``Tarskian explanation''}):

\begin{center}
\begin{tabular}{ll}
  \toprule
  $\neg S$ è vera sse & $S$ è falsa \\
  $S_1 \wedge S_2$ è vera sse & $S_1$ è vera \textbf{e} $S_2$ è vera \\
  $S_1 \vee S_2$ è vera sse & $S_1$ è vera \textbf{o} $S_2$ è vera \\
  $S_1 \to S_2$ è vera sse & $S_1$ è falsa \textbf{o} $S_2$ è vera \\
  \quad cioè è falsa sse & $S_1$ è vera \textbf{e} $S_2$ è falsa \\
  $S_1 \leftrightarrow S_2$ è vera sse & $S_1 \to S_2$ è vera \textbf{e}
      $S_2 \to S_1$ è vera \\
  \bottomrule
\end{tabular}
\end{center}

\nt{
  ``Spiegazione tarskiana'': la verità di una formula composta è definita in
  funzione della verità delle sue parti, usando nel metalinguaggio le stesse
  parole (``e'', ``o'', ``non'') che i connettivi formalizzano. La tavola di
  verità completa è nel capitolo precedente (con $1 = T$, $0 = F$).
}

\subsection{Valutazione di una formula}

Una volta fissato il significato dei simboli proposizionali, il valore di una
formula si calcola ricorsivamente con le tavole.

\ex{Valutazione ricorsiva}{
  $G \triangleq \neg P_1 \wedge (P_2 \vee P_3)$ con $P_1 = F$, $P_2 = F$,
  $P_3 = T$:
  \[
    \neg P_1 \wedge (P_2 \vee P_3) = \mathit{true} \wedge (\mathit{false} \vee
    \mathit{true}) = \mathit{true} \wedge \mathit{true} = \mathit{true}.
  \]
  Si valuta $\neg P_1$ prima di $P_1 \wedge \dots$ perché $\neg$ ha precedenza
  più alta.
}

\ex{Esercizio: XOR}{
  Tavola di verità di $G \triangleq (P \vee Q) \wedge \neg(P \wedge Q)$:
  \begin{center}
  \begin{tabular}{cc|ccc|c}
    $P$ & $Q$ & $P \vee Q$ & $P \wedge Q$ & $\neg(P \wedge Q)$ & $G$ \\ \hline
    0 & 0 & 0 & 0 & 1 & 0 \\
    0 & 1 & 1 & 0 & 1 & \textbf{1} \\
    1 & 0 & 1 & 0 & 1 & \textbf{1} \\
    1 & 1 & 1 & 1 & 0 & 0
  \end{tabular}
  \end{center}
  Annotazioni a lavagna: \textbf{ogni riga è un'interpretazione}; le righe
  $01$ e $10$ sono i \textbf{modelli} di $G$: $\{\neg P, Q\} \models G$ e
  $\{P, \neg Q\} \models G$.
}

\ex{Conseguenza logica con la tavola (lavagna)}{
  $A \wedge B,\ B \to C \models A \wedge C$? Si guardano solo le righe in cui
  \emph{tutte} le premesse sono vere.
  \begin{center}
  \begin{tabular}{ccc|cc|c}
    $A$ & $B$ & $C$ & $A \wedge B$ & $B \to C$ & $A \wedge C$ \\ \hline
    0 & 0 & 0 & 0 &   &   \\
    0 & 0 & 1 & 0 &   &   \\
    0 & 1 & 0 & 0 &   &   \\
    0 & 1 & 1 & 0 &   &   \\
    1 & 0 & 0 & 0 &   &   \\
    1 & 0 & 1 & 0 &   &   \\
    1 & 1 & 0 & 1 & 0 &   \\
    1 & 1 & 1 & 1 & 1 & \textbf{1}
  \end{tabular}
  \end{center}
  L'unica riga con entrambe le premesse vere è $111$, e lì la conclusione è
  vera: la conseguenza vale ($\checkmark$). Trucco pratico: appena una premessa
  è falsa la riga si scarta, senza calcolare il resto.
}

% =====================================================================
\section{Validità, inconsistenza, decidibilità}
% =====================================================================

\dfn{Formula valida}{
  Una formula $F$ è \textbf{valida} sse è vera in \textbf{tutte} le sue
  interpretazioni.
  \begin{itemize}
    \item Una formula valida si dice anche \textbf{tautologia}.
    \item Una formula non valida è \textbf{invalida}.
    \item Se $F$ è valida si scrive $\models F$.
  \end{itemize}
}

\ex{Legge di De Morgan}{
  $\neg(P \wedge Q) \leftrightarrow (\neg P \vee \neg Q)$ è valida:
  \begin{center}
  \begin{tabular}{cc|cc|c}
    $P$ & $Q$ & $\neg(P \wedge Q)$ & $\neg P \vee \neg Q$ &
      $\neg(P \wedge Q) \leftrightarrow (\neg P \vee \neg Q)$ \\ \hline
    1 & 1 & 0 & 0 & 1 \\
    1 & 0 & 1 & 1 & 1 \\
    0 & 1 & 1 & 1 & 1 \\
    0 & 0 & 1 & 1 & 1
  \end{tabular}
  \end{center}
}

\dfn{Formula inconsistente}{
  Una formula $F$ è \textbf{inconsistente} sse è falsa in \textbf{tutte} le
  sue interpretazioni.
  \begin{itemize}
    \item Una formula inconsistente si dice anche \textbf{insoddisfacibile}.
    \item Una formula non inconsistente è \textbf{consistente} o
          \textbf{soddisfacibile}.
    \item \textbf{Invalida è diverso da inconsistente.}
  \end{itemize}
}

\ex{}{
  $\neg\big(\neg(P \wedge Q) \leftrightarrow (\neg P \vee \neg Q)\big)$ è
  inconsistente: è la negazione della tautologia precedente, quindi la sua
  colonna è tutta $0$.
}

\mprop{Inconsistenza vs validità}{
  \begin{itemize}
    \item Una formula è valida sse la sua negazione è inconsistente (e
          viceversa).
    \item Una formula è invalida sse esiste almeno un'interpretazione in cui è
          falsa; è consistente sse esiste almeno un'interpretazione in cui è
          vera.
    \item Una formula inconsistente è invalida, \textbf{ma non vale il
          viceversa}.
    \item Una formula valida è consistente, \textbf{ma non vale il viceversa}.
  \end{itemize}
}

\ex{}{
  $G \triangleq P \vee Q$ è \textbf{invalida} (falsa quando $P$ e $Q$ sono
  falsi) ma \textbf{non inconsistente} (vera in tutti gli altri casi). Detto
  al contrario: è \textbf{consistente} ma \textbf{non valida}.
}

\nt{
  Schema riassuntivo: ogni formula sta in esattamente una di tre classi.
  \begin{center}
  \begin{tabular}{lccc}
    \toprule
    & sempre vera & a volte vera, a volte falsa & sempre falsa \\ \midrule
    valida (tautologia) & $\checkmark$ & & \\
    soddisfacibile / consistente & $\checkmark$ & $\checkmark$ & \\
    invalida & & $\checkmark$ & $\checkmark$ \\
    insoddisfacibile / inconsistente & & & $\checkmark$ \\
    \bottomrule
  \end{tabular}
  \end{center}
}

\subsection{Le tre domande e la decidibilità}

Sia $AX$ un insieme di assiomi che descrive un ``mondo'' e $F$ una formula.
Ci interessano tre domande:
\begin{enumerate}
  \item $AX \models F$? Cioè $F$ è conseguenza logica di $AX$?
  \item $\models F$? Cioè $F$ è valida?
  \item $F$ è soddisfacibile? Cioè esiste $I$ tale che $I \models F$?
\end{enumerate}
Esempio per la 3: $F$ descrive un guasto in un circuito elettrico con
variabili $X_1, \dots, X_n$; sapere se è soddisfacibile vuol dire sapere se il
guasto può verificarsi. È importante sapere se a queste domande si può
\emph{sempre} rispondere sì o no, cioè se i problemi sono \textbf{decidibili}.

\thm{La logica proposizionale è decidibile}{
  Esiste un metodo che \textbf{termina sempre} per decidere se una formula è
  valida:
  \begin{enumerate}
    \item si enumerano tutte le interpretazioni possibili;
    \item per ciascuna si valuta la formula.
  \end{enumerate}
  Le interpretazioni di una formula sono in numero finito ($2^n$).
}

\begin{itemize}
  \item La decidibilità è una proprietà molto forte e desiderabile per un
        sistema logico.
  \item C'è un \textbf{trade-off tra potere espressivo e decidibilità} (la
        logica del primo ordine è più espressiva ma non decidibile).
\end{itemize}

\subsection{Decidere la soddisfacibilità: Davis--Putnam}

Gli algoritmi più veloci noti per la soddisfacibilità proposizionale si basano
sull'\textbf{algoritmo di Davis--Putnam}. Lavora su una lista di clausole
(formula in CNF, vedi sotto). Una \textbf{unit clause} è una clausola formata
da un solo letterale.

\begin{lstlisting}[mathescape=true]
function Satisfiable(clause list S) returns boolean
    /* unit propagation */
    repeat
        for each unit clause L in S do
            delete from S every clause containing L
            delete $\neg$L from every clause of S in which it occurs
        end for
        if S is empty then return TRUE
        else if null clause is in S then return FALSE
    until no further changes result
    /* splitting */
    choose a literal L occurring in S
    if Satisfiable(S $\cup$ {L}) then return TRUE
    else if Satisfiable(S $\cup$ {$\neg$L}) then return TRUE
    else return FALSE
end function
\end{lstlisting}

\nt{
  \textbf{Come leggerlo.}
  \begin{itemize}
    \item \emph{Unit propagation}: se c'è la clausola unitaria $L$, $L$ deve
          essere vero. Allora ogni clausola che contiene $L$ è già soddisfatta
          (si cancella) e in ogni altra clausola $\neg L$ è falso (si toglie il
          letterale).
    \item $S$ vuota: non restano vincoli, \textbf{soddisfacibile}. Clausola
          vuota (null clause): una disgiunzione di zero letterali è falsa,
          \textbf{insoddisfacibile} su questo ramo.
    \item \emph{Splitting}: si sceglie un letterale e si provano i due casi
          ($L$ vero, $L$ falso), aggiungendolo come clausola unitaria.
          È una ricerca con backtracking: caso peggiore esponenziale (SAT è
          NP-completo), ma in pratica molto più veloce della tavola di verità.
  \end{itemize}
}

\ex{Davis--Putnam a mano}{
  $S = \{\,P \vee Q,\ \neg P \vee R,\ \neg R,\ \neg Q \vee R\,\}$.
  \begin{enumerate}
    \item Unit clause $\neg R$: cancello $\neg R$; tolgo $R$ da
          $\neg P \vee R$ e $\neg Q \vee R$. Resta
          $\{P \vee Q,\ \neg P,\ \neg Q\}$.
    \item Unit clause $\neg P$: cancello $\neg P$; tolgo $P$ da $P \vee Q$.
          Resta $\{Q,\ \neg Q\}$.
    \item Unit clause $Q$: cancello $Q$; tolgo $\neg Q$ da $\neg Q$, che
          diventa la \textbf{clausola vuota}. Risposta: FALSE,
          $S$ è insoddisfacibile.
  \end{enumerate}
}

% =====================================================================
\section{Equivalenza logica e regole di riscrittura}
% =====================================================================

\dfn{Equivalenza logica}{
  Due formule $F$ e $G$ sono \textbf{logicamente equivalenti}, $F \equiv G$,
  sse hanno lo stesso valore di verità in ogni interpretazione di $F$ e $G$.
}

\thm{Regole di equivalenza utili}{
  \begin{center}
  \begin{tabular}{rcll}
    $P \wedge Q$ & $\equiv$ & $Q \wedge P$ & commutatività di $\wedge$ \\
    $P \vee Q$ & $\equiv$ & $Q \vee P$ & commutatività di $\vee$ \\
    $(P \wedge Q) \wedge R$ & $\equiv$ & $P \wedge (Q \wedge R)$ & associatività di $\wedge$ \\
    $(P \vee Q) \vee R$ & $\equiv$ & $P \vee (Q \vee R)$ & associatività di $\vee$ \\
    $\neg(\neg P)$ & $\equiv$ & $P$ & eliminazione della doppia negazione \\
    $P \to Q$ & $\equiv$ & $\neg Q \to \neg P$ & contrapposizione \\
    $P \to Q$ & $\equiv$ & $\neg P \vee Q$ & eliminazione dell'implicazione \\
    $P \leftrightarrow Q$ & $\equiv$ & $(P \to Q) \wedge (Q \to P)$ & eliminazione del bicondizionale \\
    $\neg(P \wedge Q)$ & $\equiv$ & $\neg P \vee \neg Q$ & De Morgan \\
    $\neg(P \vee Q)$ & $\equiv$ & $\neg P \wedge \neg Q$ & De Morgan \\
    $P \wedge (Q \vee R)$ & $\equiv$ & $(P \wedge Q) \vee (P \wedge R)$ & distributività di $\wedge$ su $\vee$ \\
    $P \vee (Q \wedge R)$ & $\equiv$ & $(P \vee Q) \wedge (P \vee R)$ & distributività di $\vee$ su $\wedge$ \\
  \end{tabular}
  \end{center}
}

\thm{Regole con $\top$ e $\bot$ (aggiunte a lavagna)}{
  \begin{center}
  \begin{tabular}{rclcrcll}
    $P \wedge \top$ & $\equiv$ & $P$ & \qquad & $P \vee \bot$ & $\equiv$ & $P$ & identità \\
    $P \wedge \bot$ & $\equiv$ & $\bot$ & & $P \vee \top$ & $\equiv$ & $\top$ & elemento assorbente \\
    $(P \wedge Q) \vee Q$ & $\equiv$ & $Q$ & & $(P \vee Q) \wedge Q$ & $\equiv$ & $Q$ & assorbimento \\
    $\neg\bot$ & $\equiv$ & $\top$ & & $\neg\top$ & $\equiv$ & $\bot$ & \\
  \end{tabular}
  \end{center}
}

\nt{
  Per l'assorbimento conviene pensare all'ordine ``$0 \le 1$'':
  $P \wedge Q$ è ``più falso'' di $Q$ (è il minimo), quindi
  $(P \wedge Q) \vee Q = \max(\min(P,Q), Q) = Q$; dualmente $P \vee Q$ è ``più
  vero'' di $Q$, quindi $(P \vee Q) \wedge Q = Q$. È il disegno della retta
  fatto a lavagna.
}

% =====================================================================
\section{Forme normali}
% =====================================================================

Le forme normali sono modi \emph{standard} di scrivere le formule. Le due
principali: \textbf{CNF} (Conjunctive Normal Form) e \textbf{DNF} (Disjunctive
Normal Form).

\dfn{Letterale}{
  Un \textbf{letterale} è un atomo (a mano: cioè una variabile) o la negazione
  di un atomo. Esempi: $A$, $\neg B$, $C$.
}

\dfn{Negation Normal Form (NNF)}{
  Una formula è in \textbf{NNF} sse le negazioni compaiono \textbf{solo
  davanti agli atomi}.
}

\dfn{Conjunctive Normal Form (CNF)}{
  Una formula $F$ è in \textbf{CNF} sse è in NNF e ha la forma
  \[ F \triangleq F_1 \wedge F_2 \wedge \dots \wedge F_n \]
  dove ogni $F_i$ è una \textbf{disgiunzione di letterali}.
  \begin{itemize}
    \item Ogni $F_i$ si chiama \textbf{clausola}: una clausola è una formula
          $L_1 \vee \dots \vee L_n$ con ogni $L_i$ letterale.
    \item La CNF si chiama anche \textbf{forma a clausole} (Clausal Form).
  \end{itemize}
}

\nt{
  A lavagna: l'albero di una CNF ha \textbf{un $\wedge$ in radice, sotto solo
  $\vee$, sotto solo letterali}, ad esempio
  $(A \vee \neg B \vee C) \wedge (D \vee A) \wedge A$. Per la DNF è il
  contrario: $\vee$ in radice, poi $\wedge$, poi letterali.
}

\ex{}{
  \begin{itemize}
    \item $G \triangleq (\neg P \vee Q) \wedge (\neg P \vee R)$ è in CNF. Si
          può scrivere come \textbf{insieme di clausole} $\{C_1, C_2\}$ con
          $C_1 = \neg P \vee Q$ e $C_2 = \neg P \vee R$.
    \item $G \triangleq \neg(P \vee Q) \wedge (\neg P \vee R)$ \textbf{non} è
          in CNF: la negazione compare davanti a una formula e non solo davanti
          ad atomi.
  \end{itemize}
}

\dfn{Disjunctive Normal Form (DNF)}{
  Una formula $F$ è in \textbf{DNF} sse è in NNF e ha la forma
  \[ F \triangleq F_1 \vee F_2 \vee \dots \vee F_n \]
  dove ogni $F_i$ è una \textbf{congiunzione di letterali}.
  Esempio: $(\neg P \wedge R) \vee (Q \wedge \neg P) \vee (Q \wedge P)$.
}

\thm{Esistenza delle forme normali}{
  Ogni formula può essere trasformata in una forma normale (CNF o DNF)
  equivalente usando le regole di equivalenza viste sopra.
}

\subsection{Algoritmo di trasformazione}

\mprop{Trasformazione in CNF in 5 passi (lavagna)}{
  \begin{enumerate}
    \item \textbf{Rimuovere} $\to$ e $\leftrightarrow$
          \hfill ($F \to G \equiv \neg F \vee G$, ecc.);
    \item \textbf{spingere le negazioni verso il basso} (verso gli atomi)
          \hfill (De Morgan);
    \item \textbf{rimuovere} $\neg\neg$ \hfill (doppia negazione);
    \item \textbf{distribuire $\vee$ su $\wedge$}
          \hfill ($P \vee (Q \wedge R) \equiv (P \vee Q) \wedge (P \vee R)$);
    \item \textbf{semplificare} $\bot$ / $\top$.
  \end{enumerate}
  Per la \textbf{DNF} il passo 4 diventa: distribuire $\wedge$ su $\vee$.
}

\ex{CNF svolta a lavagna}{
  \begin{align*}
    & (P \to Q) \wedge \neg(\neg R \wedge (B \vee A)) \\
    \equiv\ & (\neg P \vee Q) \wedge (\neg\neg R \vee \neg(B \vee A))
       && \text{1: via $\to$; 2: De Morgan} \\
    \equiv\ & (\neg P \vee Q) \wedge (\neg\neg R \vee (\neg B \wedge \neg A))
       && \text{2: De Morgan} \\
    \equiv\ & (\neg P \vee Q) \wedge (R \vee (\neg B \wedge \neg A))
       && \text{3: via $\neg\neg$} \\
    \equiv\ & (\neg P \vee Q) \wedge (R \vee \neg B) \wedge (R \vee \neg A)
       && \text{4: distribuzione} \quad \textbf{CNF}
  \end{align*}
}

\ex{Trasformazione in CNF (slide)}{
  \begin{align*}
    & (P \to Q) \wedge (R \vee (B \wedge A)) \\
    \equiv\ & (\neg P \vee Q) \wedge (R \vee (B \wedge A)) \\
    \equiv\ & (\neg P \vee Q) \wedge ((R \vee B) \wedge (R \vee A)) \\
    \equiv\ & (\neg P \vee Q) \wedge (R \vee B) \wedge (R \vee A)
  \end{align*}
}

\ex{Altro esempio (slide): semplificazione di $\top$ e $\bot$}{
  \begin{align*}
    & (\neg P \vee Q) \wedge (R \vee (A \wedge \neg A)) \wedge
      (R \vee (A \vee \neg A)) \\
    \equiv\ & (\neg P \vee Q) \wedge (R \vee \bot) \wedge (R \vee \top) \\
    \equiv\ & (\neg P \vee Q) \wedge (R \vee \bot) \wedge \top \\
    \equiv\ & (\neg P \vee Q) \wedge (R \vee \bot) \\
    \equiv\ & (\neg P \vee Q) \wedge R
  \end{align*}
}

\nt{
  \textbf{Attenzione, probabile refuso nella slide.} Sulla slide la seconda
  clausola è scritta $R \vee (A \to \neg A)$ e viene semplificata in
  $R \vee \mathit{False}$. Ma $A \to \neg A \equiv \neg A \vee \neg A \equiv
  \neg A$, che \textbf{non} è equivalente a falso (è vera quando $A$ è falso).
  Il passaggio è corretto solo con una contraddizione, cioè
  $A \wedge \neg A \equiv \bot$: sopra ho scritto questa versione.
}

\ex{Trasformazione di formule: XOR}{
  Dimostrare, trasformando le formule, che
  \[
    (P \vee Q) \wedge \neg(P \wedge Q) \;\equiv\;
    (P \vee Q) \wedge (\neg P \vee \neg Q) \;\equiv\;
    (\neg P \wedge Q) \vee (P \wedge \neg Q).
  \]
  \textbf{Prima equivalenza} (slide): De Morgan sulla seconda parte dà
  direttamente $(P \vee Q) \wedge (\neg P \vee \neg Q)$, che è la CNF.

  \textbf{Seconda equivalenza} (``try to prove''): distribuendo $\wedge$ su
  $\vee$,
  \begin{align*}
    (P \vee Q) \wedge (\neg P \vee \neg Q)
    &\equiv (P \wedge \neg P) \vee (P \wedge \neg Q) \vee (Q \wedge \neg P)
       \vee (Q \wedge \neg Q) \\
    &\equiv \bot \vee (P \wedge \neg Q) \vee (\neg P \wedge Q) \vee \bot \\
    &\equiv (\neg P \wedge Q) \vee (P \wedge \neg Q),
  \end{align*}
  che è la DNF. (Altri esempi: Chang--Lee, esempi 2.8 e 2.9.)
}

\ex{Esame 06/07/2021, es.\ 2: mettere in DNF}{
  \begin{align*}
    & \neg(A \to B \vee C) \;\vee\; ((B \vee A) \to C \wedge D) \\
    \equiv\ & \neg(\neg A \vee B \vee C) \;\vee\; \neg(B \vee A) \vee (C \wedge D)
       && \text{via le due $\to$} \\
    \equiv\ & (A \wedge \neg B \wedge \neg C) \;\vee\; (\neg B \wedge \neg A)
       \;\vee\; (C \wedge D)
       && \text{De Morgan, $\neg\neg$}
  \end{align*}
  Attenzione alle precedenze: $A \to B \vee C$ è $A \to (B \vee C)$ e
  $(B \vee A) \to C \wedge D$ è $(B \vee A) \to (C \wedge D)$.
}

% =====================================================================
\section{Conseguenza logica, deduzione e refutazione}
% =====================================================================

\dfn{Conseguenza logica}{
  Dati un insieme di formule $\{F_1, \dots, F_n\}$ e una formula $G$, $G$ è
  \textbf{conseguenza logica} di $F_1, \dots, F_n$ sse in ogni interpretazione
  $I$ in cui $F_1 \wedge \dots \wedge F_n$ è vera anche $G$ è vera. Si scrive
  \[ F_1 \wedge \dots \wedge F_n \models G. \]
  \begin{itemize}
    \item $F_1, \dots, F_n$ si chiamano \textbf{assiomi} o \textbf{premesse}
          di $G$.
    \item $F \equiv Q$ sse $F \models Q$ e $Q \models F$.
  \end{itemize}
  Esempio: $S \to C$, $C \to F$, $S$ sono premesse per $F$.
}

\thm{Teorema di deduzione}{
  Dati $\{F_1, \dots, F_n\}$ e $G$:
  \[
    (F_1 \wedge \dots \wedge F_n) \models G
    \quad\text{sse}\quad
    \models (F_1 \wedge \dots \wedge F_n) \to G.
  \]
}

\pf{Schema della dimostrazione}{
  ($\Rightarrow$) In ogni interpretazione $I$ in cui $F_1 \wedge \dots \wedge
  F_n$ è vera, $G$ è vera, quindi $I \models (F_1 \wedge \dots \wedge F_n) \to
  G$. In ogni interpretazione $I'$ in cui $F_1 \wedge \dots \wedge F_n$ è
  falsa, l'implicazione è vera (premessa falsa), quindi $I' \models (F_1
  \wedge \dots \wedge F_n) \to G$. Dunque l'implicazione è vera in tutte le
  interpretazioni: è valida.

  ($\Leftarrow$) Se l'implicazione è valida, in ogni interpretazione in cui
  $F_1 \wedge \dots \wedge F_n$ è vera deve essere vera anche $G$ (altrimenti
  l'implicazione sarebbe falsa). Quindi $(F_1 \wedge \dots \wedge F_n) \models
  G$.
}

\thm{Dimostrazione per refutazione}{
  Dati $\{F_1, \dots, F_n\}$ e $G$:
  \[
    (F_1 \wedge \dots \wedge F_n) \models G
    \quad\text{sse}\quad
    F_1 \wedge \dots \wedge F_n \wedge \neg G \ \text{è inconsistente.}
  \]
}

\pf{Schema della dimostrazione}{
  $(F_1 \wedge \dots \wedge F_n) \models G$ vale sse in ogni interpretazione
  in cui $F_1 \wedge \dots \wedge F_n$ è vera anche $G$ è vera. Questo vale
  sse non c'è nessuna interpretazione in cui $F_1 \wedge \dots \wedge F_n$ è
  vera e $G$ è falsa, cioè esattamente quando $F_1 \wedge \dots \wedge F_n
  \wedge \neg G$ è falsa in ogni interpretazione: è inconsistente.
}

\mprop{Discussione}{
  I due teoremi mostrano che:
  \begin{itemize}
    \item si può dimostrare una conseguenza logica dimostrando la
          \textbf{validità} di una formula;
    \item si può dimostrare una conseguenza logica \textbf{refutando} una
          formula, cioè dimostrando che è inconsistente.
  \end{itemize}
  \textbf{Non abbiamo usato nessuna proprietà specifica della logica
  proposizionale}: gli stessi teoremi valgono (con gli stessi argomenti) per la
  logica del primo ordine, ed è su questo che si basa la risoluzione.
}

\ex{Modus ponens nei tre modi}{
  Vogliamo mostrare $(P \to Q) \wedge P \models Q$.
  \begin{center}
  \begin{tabular}{cc|c|c|c}
    $P$ & $Q$ & $(P \to Q) \wedge P$ & $((P \to Q) \wedge P) \to Q$ &
      $(P \to Q) \wedge P \wedge \neg Q$ \\ \hline
    0 & 0 & 0 & 1 & 0 \\
    0 & 1 & 0 & 1 & 0 \\
    1 & 0 & 0 & 1 & 0 \\
    1 & 1 & 1 & 1 & 0
  \end{tabular}
  \end{center}
  \begin{itemize}
    \item \textbf{Con la definizione}: nell'unica interpretazione in cui
          $(P \to Q) \wedge P$ è vera (riga $11$), $Q$ è vera.
    \item \textbf{Con il teorema di deduzione}: $((P \to Q) \wedge P) \to Q$ è
          valida (colonna tutta $1$).
    \item \textbf{Per refutazione}: $(P \to Q) \wedge P \wedge \neg Q$ è
          inconsistente (colonna tutta $0$).
  \end{itemize}
}

\ex{Esercizio (Chang--Lee 2.11): modus tollens nei tre modi}{
  $F_1 \triangleq P \to Q$, $F_2 \triangleq \neg Q$, $G \triangleq \neg P$.
  Mostrare $F_1 \wedge F_2 \models G$.
  \begin{center}
  \begin{tabular}{cc|cc|c|c|c}
    $P$ & $Q$ & $F_1$ & $F_2$ & $G$ & $(F_1 \wedge F_2) \to G$ &
      $F_1 \wedge F_2 \wedge \neg G$ \\ \hline
    0 & 0 & 1 & 1 & 1 & 1 & 0 \\
    0 & 1 & 1 & 0 & 1 & 1 & 0 \\
    1 & 0 & 0 & 1 & 0 & 1 & 0 \\
    1 & 1 & 1 & 0 & 0 & 1 & 0
  \end{tabular}
  \end{center}
  (1) L'unica riga con $F_1$ e $F_2$ vere è $00$, e lì $G$ è vera. (2)
  $(F_1 \wedge F_2) \to G$ è valida. (3) $F_1 \wedge F_2 \wedge \neg G$ è
  inconsistente. In alternativa al punto (3), per trasformazione:
  $(\neg P \vee Q) \wedge \neg Q \wedge P \equiv ((\neg P \wedge \neg Q) \vee
  (Q \wedge \neg Q)) \wedge P \equiv \neg P \wedge \neg Q \wedge P \equiv
  \bot$.
}

\ex{Esercizio (Chang--Lee 2.12): lo sciopero}{
  ``Dato che, se il congresso rifiuta di approvare nuove leggi, allora lo
  sciopero non finirà a meno che non duri più di un anno o il presidente
  dell'azienda si dimetta: lo sciopero finirà se il congresso rifiuta di agire
  e lo sciopero è appena cominciato?''

  Atomi: $P$ = il congresso rifiuta di agire; $Q$ = lo sciopero finisce;
  $R$ = il presidente si dimette; $S$ = lo sciopero dura più di un anno.
  \begin{itemize}
    \item $F_1 \triangleq P \to (\neg Q \vee R \vee S)$ \quad (``non $Q$ a
          meno che $R$ o $S$'' = $Q \to (R \vee S)$);
    \item $F_2 \triangleq P$; \quad $F_3 \triangleq \neg S$ (è appena
          cominciato).
  \end{itemize}
  Da $F_1, F_2$ per modus ponens: $\neg Q \vee R \vee S$; con $F_3$:
  $\neg Q \vee R$, cioè $Q \to R$. Con la formulazione della slide (``\emph{or}
  the president resigns'') \textbf{non segue né $Q$ né $\neg Q$}:
  l'interpretazione $\{P, Q, R, \neg S\}$ rende vere le premesse con $Q$ vero,
  $\{P, \neg Q, \neg R, \neg S\}$ con $Q$ falso. Quello che si può dire è solo
  che lo sciopero finisce \emph{solo se} il presidente si dimette.

  Se invece la condizione è una congiunzione (``dura più di un anno \emph{e}
  il presidente si dimette''), $F_1 \triangleq P \to (\neg Q \vee (R \wedge
  S))$ e allora $F_1 \wedge F_2 \wedge F_3 \models \neg Q$: lo sciopero
  \textbf{non} finirà. Per refutazione: $F_1 \wedge P \wedge \neg S \wedge Q$
  forza $R \wedge S$ vero, contro $\neg S$.
}

\ex{Esercizio: argomenti validi e non validi}{
  \begin{enumerate}
    \item ``Se sono ricco allora sono felice. Sono felice. Quindi sono
          ricco.'' \ $w \to h,\ h \models w$? \textbf{Non valido}
          (affermazione del conseguente): controesempio $\{\neg w, h\}$.
    \item ``Se John beve birra ha almeno 18 anni. John non beve birra. Quindi
          non ha ancora 18 anni.'' \ $b \to a,\ \neg b \models \neg a$?
          \textbf{Non valido} (negazione dell'antecedente): controesempio
          $\{\neg b, a\}$.
    \item ``Se le ragazze sono bionde, sono popolari. Le brutte sono
          impopolari. Le intellettuali sono brutte. Quindi le bionde non sono
          intellettuali.'' \ $b \to p,\ u \to \neg p,\ i \to u \models b \to
          \neg i$? \textbf{Valido}: $b \to p$, poi $p \to \neg u$
          (contrapposizione della seconda), poi $\neg u \to \neg i$
          (contrapposizione della terza); per transitività $b \to \neg i$.
    \item ``Se studio non vengo bocciato. Se non gioco troppo a carte, studio.
          Sono stato bocciato. Quindi ho giocato troppo a carte.''
          \ $s \to \neg f,\ \neg c \to s,\ f \models c$? \textbf{Valido}:
          da $f$ e $s \to \neg f$ (modus tollens) $\neg s$; da $\neg s$ e
          $\neg c \to s$ (modus tollens) $c$.
  \end{enumerate}
}

\ex{Esercizio: le lettere di Lewis Carroll}{
  Atomi: $p$ = datata; $q$ = su carta blu; $r$ = in inchiostro nero; $s$ = in
  terza persona; $t$ = archiviata; $u$ = la so leggere; $v$ = su un solo
  foglio; $w$ = cancellata; $x$ = scritta da Brown; $y$ = comincia con ``Dear
  Sir''. Le nove premesse:
  \begin{multicols}{3}
  \begin{enumerate}
    \item $p \to q$
    \item $\neg s \to \neg r$
    \item $u \to \neg t$
    \item $v \to p$
    \item $\neg w \to r$
    \item $x \to y$
    \item $q \to t$
    \item $\neg v \to \neg w$
    \item $y \to \neg s$
  \end{enumerate}
  \end{multicols}
  Tesi: $x \to \neg u$ (``non so leggere nessuna lettera di Brown'').

  \textbf{Soluzione.} Si costruisce una catena di implicazioni, usando la
  contrapposizione dove serve:
  \[
    x \xrightarrow{6} y \xrightarrow{9} \neg s \xrightarrow{2} \neg r
      \xrightarrow{5^c} w \xrightarrow{8^c} v \xrightarrow{4} p
      \xrightarrow{1} q \xrightarrow{7} t \xrightarrow{3^c} \neg u
  \]
  dove $5^c$ è $\neg r \to w$, $8^c$ è $w \to v$, $3^c$ è $t \to \neg u$.
  Per transitività dell'implicazione $x \to \neg u$ è conseguenza logica delle
  premesse: l'argomento è valido. (Con la tavola di verità servirebbero
  $2^{10} = 1024$ righe: ecco perché serve un metodo \emph{sintattico}.)
}

\ex{Esame 12/01/2023, es.\ 2}{
  \begin{itemize}
    \item $(A \to B) \wedge (A \to C) \models A \wedge B \to C$? \textbf{Sì}
          ($\checkmark$). Le righe in cui la premessa è vera sono
          $000, 001, 010, 011, 111$; nelle prime quattro $A \wedge B$ è falso
          (implicazione vera), nell'ultima $C$ è vero.
    \item $A \vee B \vee \neg C \models A \wedge B \to C$? \textbf{No}
          ($\times$). Controesempio: riga $A = 1, B = 1, C = 0$: la premessa è
          vera, $A \wedge B = 1$ e $C = 0$, quindi la conclusione è falsa.
  \end{itemize}
}

\ex{Esame 12/01/2023, es.\ 3: un insieme non contraddittorio}{
  $M \to C,\ \neg M \to T,\ C \vee T \vdash \bot$? \textbf{No}: l'insieme
  \textbf{non è contraddittorio}. Dalla tavola, le righe $(M,C,T)$ che rendono
  vere tutte e tre le formule sono
  \[ 001,\quad 011,\quad 110,\quad 111. \]
  Basta un modello per concludere che l'insieme è soddisfacibile, quindi (per
  la correttezza, vedi sotto) da esso non si deriva $\bot$.
}

% =====================================================================
\section{Semantica II: proof theory}
% =====================================================================

\begin{itemize}
  \item Ci interessano le \textbf{conseguenze logiche}: permettono di ricavare
        conclusioni sulla verità in mondi specifici descritti da assiomi
        specifici.
  \item Abbiamo visto un modo \textbf{semantico} di dimostrarle, basato sulla
        nozione di modello: le tavole di verità.
  \item Esiste un modo di dimostrarle usando la \textbf{sintassi}, cioè
        manipolando i simboli delle formule? Questo permetterebbe di usare la
        \emph{computazione} per fare \textbf{deduzione (o inferenza)} in
        automatico.
  \item La risposta è sì, e la dà la \textbf{proof theory}.
\end{itemize}

Esistono molti sistemi, per molte logiche diverse, fatti di \textbf{regole}
che permettono di dedurre (derivare) \textbf{conclusioni da premesse}.
Naturalmente vogliamo che le regole siano \textbf{corrette} rispetto alla
semantica delle tavole di verità. Vediamo un caso specifico: la
\textbf{deduzione naturale} per (a mano: un \textbf{sottoinsieme} della)
logica proposizionale \textbf{classica}.

\nt{
  \textbf{Gli alberi (lavagna).} Un albero vero ha le radici in basso e le
  foglie in alto; gli alberi dell'informatica hanno la radice in alto. Le
  \textbf{derivazioni} sono disegnate come l'albero vero: \textbf{foglie in
  alto} (le ipotesi), \textbf{radice in basso} (la conclusione). Si possono
  leggere \emph{top-down} (dalle ipotesi alla conclusione) o costruire
  \emph{bottom-up} (dalla conclusione verso le ipotesi), che è il modo in cui
  di solito si cerca una dimostrazione.
}

\subsection{Deduzione naturale}

\begin{itemize}
  \item Un insieme di \textbf{regole} per derivare conclusioni da premesse.
  \item Le regole \textbf{introducono} o \textbf{eliminano} connettivi.
  \item Le regole si compongono a formare \textbf{derivazioni}.
  \item È un sistema \textbf{corretto e completo} (vedi oltre).
\end{itemize}

\mprop{Connettivi considerati}{
  Consideriamo solo $\wedge$, $\to$ e $\bot$. Il motivo è che $\vee$ ha un
  significato molto diverso nell'approccio costruttivo e in quello non
  costruttivo. Non è una limitazione: questo insieme di connettivi è
  \textbf{funzionalmente completo}.
}

\thm{Gli altri connettivi come abbreviazioni (lavagna)}{
  Teniamo $\wedge, \to, \bot$ e definiamo:
  \begin{align*}
    \neg F &\equiv F \to \bot \\
    \top &\equiv \neg\bot \equiv \bot \to \bot \\
    F \leftrightarrow G &\equiv (F \to G) \wedge (G \to F) \\
    F \vee G &\equiv \neg\neg(F \vee G) \equiv \neg(\neg F \wedge \neg G)
      \equiv \big((F \to \bot) \wedge (G \to \bot)\big) \to \bot
  \end{align*}
  L'ultima riga vale \textbf{solo in logica classica}.
}

\nt{
  Nella logica intuizionista (costruttiva) dimostrare $F \vee G$ significa
  saper dire \emph{quale} dei due vale; $\neg(\neg F \wedge \neg G)$ dice solo
  che non possono essere entrambi falsi, senza dire quale è vero. Per questo
  la traduzione di $\vee$ è accettabile solo classicamente.
}

\dfn{Regole della deduzione naturale}{
  \begin{center}
  \begin{tabular}{c@{\qquad\qquad}c}
    \textsc{introduzione} & \textsc{eliminazione} \\[2ex]
    \AxiomC{$\varphi$}
    \AxiomC{$\psi$}
    \RightLabel{$\wedge I$}
    \BinaryInfC{$\varphi \wedge \psi$}
    \DisplayProof
    &
    \AxiomC{$\varphi \wedge \psi$}
    \RightLabel{$\wedge E_1$}
    \UnaryInfC{$\varphi$}
    \DisplayProof
    \quad
    \AxiomC{$\varphi \wedge \psi$}
    \RightLabel{$\wedge E_2$}
    \UnaryInfC{$\psi$}
    \DisplayProof
    \\[5ex]
    \AxiomC{$[\varphi]$}
    \noLine
    \UnaryInfC{$\vdots$}
    \noLine
    \UnaryInfC{$\psi$}
    \RightLabel{$\to I$}
    \UnaryInfC{$\varphi \to \psi$}
    \DisplayProof
    &
    \AxiomC{$\varphi$}
    \AxiomC{$\varphi \to \psi$}
    \RightLabel{$\to E$}
    \BinaryInfC{$\psi$}
    \DisplayProof
  \end{tabular}
  \end{center}
  Le parentesi quadre $[\ ]$ indicano lo \textbf{scarico} (cancellazione)
  dell'ipotesi.

  Per $\bot$ ci sono due regole; entrambe eliminano $\bot$ e introducono una
  formula:
  \begin{center}
  \begin{tabular}{c@{\qquad\qquad}c}
    \emph{Ex falso sequitur quodlibet} & \emph{Reductio ad absurdum} \\[2ex]
    \AxiomC{$\bot$}
    \RightLabel{$\bot_E$}
    \UnaryInfC{$\varphi$}
    \DisplayProof
    &
    \AxiomC{$[\neg\varphi]$}
    \noLine
    \UnaryInfC{$\vdots$}
    \noLine
    \UnaryInfC{$\bot$}
    \RightLabel{RAA}
    \UnaryInfC{$\varphi$}
    \DisplayProof
  \end{tabular}
  \end{center}
  $\neg\varphi$ è usato come abbreviazione di $\varphi \to \bot$.
  \textbf{RAA vale solo per la logica classica.}
}

\nt{
  \textbf{Come leggere le regole.}
  \begin{itemize}
    \item $\wedge I$: se ho $\varphi$ e ho $\psi$, ho $\varphi \wedge \psi$.
          $\wedge E$: da $\varphi \wedge \psi$ ricavo ciascuna delle due.
    \item $\to E$ è il \emph{modus ponens}.
    \item $\to I$: per dimostrare $\varphi \to \psi$ \emph{assumo} $\varphi$,
          derivo $\psi$, e poi \textbf{scarico} l'assunzione: la conclusione
          $\varphi \to \psi$ non dipende più da $\varphi$. Si possono scaricare
          zero, una o più occorrenze di $\varphi$ (anche \emph{nessuna}: da
          $\psi$ ricavo $\varphi \to \psi$ per qualunque $\varphi$).
    \item $\bot_E$: dall'assurdo segue qualsiasi cosa.
    \item RAA: se assumere $\neg\varphi$ porta all'assurdo, allora $\varphi$.
          Attenzione alla differenza con $\to I$: se da $\varphi$ derivo
          $\bot$, con $\to I$ ottengo $\varphi \to \bot = \neg\varphi$
          (costruttivo). Con RAA da $\neg\varphi \vdash \bot$ ottengo
          $\varphi$ e non solo $\neg\neg\varphi$: è qui che entra la logica
          classica.
    \item Gli indici $[\varphi]^1$ e $\to I_1$ dicono \emph{quale} regola
          scarica \emph{quale} ipotesi.
  \end{itemize}
}

\mprop{Regole invertibili (lavagna)}{
  \begin{center}
  \begin{tabular}{ll}
    \toprule
    \textbf{invertibili} & $\wedge I$, \quad $\to I$, \quad RAA \\
    \textbf{non invertibili} & $\wedge E$, \quad $\to E$, \quad $\bot_E$ \\
    \bottomrule
  \end{tabular}
  \end{center}
}

\nt{
  (Approfondimento.) Una regola è \textbf{invertibile} se, quando la
  conclusione è derivabile, lo sono anche le premesse: applicarla
  \emph{bottom-up} non fa mai perdere la dimostrabilità, quindi si applica
  ``senza pensare''.
  \begin{itemize}
    \item $\wedge I$: se $\Gamma \vdash \varphi \wedge \psi$ allora
          $\Gamma \vdash \varphi$ e $\Gamma \vdash \psi$.
    \item $\to I$: se $\Gamma \vdash \varphi \to \psi$ allora
          $\Gamma, \varphi \vdash \psi$.
    \item RAA: se $\Gamma \vdash \varphi$ allora $\Gamma, \neg\varphi \vdash
          \bot$.
    \item $\wedge E$ no: per derivare $\varphi$ devo \emph{indovinare} un
          $\psi$ tale che $\varphi \wedge \psi$ sia derivabile. $\to E$ no:
          devo indovinare $\varphi$. $\bot_E$ no: $\varphi$ può essere
          derivabile senza che lo sia $\bot$.
  \end{itemize}
  \textbf{Strategia pratica}: partire dal basso applicando $\to I$ e
  $\wedge I$ finché si può (scrivendo in alto le ipotesi disponibili), poi
  combinare le ipotesi con le regole di eliminazione; se non si chiude,
  provare RAA.
}

\subsection{Esempi di derivazioni}

\ex{I: $\vdash \varphi \wedge \psi \to \psi \wedge \varphi$}{
  \begin{prooftree}
    \AxiomC{$[\varphi \wedge \psi]^1$}
    \RightLabel{$\wedge E$}
    \UnaryInfC{$\psi$}
    \AxiomC{$[\varphi \wedge \psi]^1$}
    \RightLabel{$\wedge E$}
    \UnaryInfC{$\varphi$}
    \RightLabel{$\wedge I$}
    \BinaryInfC{$\psi \wedge \varphi$}
    \RightLabel{$\to I_1$}
    \UnaryInfC{$\varphi \wedge \psi \to \psi \wedge \varphi$}
  \end{prooftree}
  A lavagna la versione con ipotesi \emph{non} scaricata, cioè
  $A \wedge B \vdash B \wedge A$: è lo stesso albero senza l'ultimo passo.
  \begin{prooftree}
    \AxiomC{$A \wedge B$}
    \RightLabel{$\wedge E_2$}
    \UnaryInfC{$B$}
    \AxiomC{$A \wedge B$}
    \RightLabel{$\wedge E_1$}
    \UnaryInfC{$A$}
    \RightLabel{$\wedge I$}
    \BinaryInfC{$B \wedge A$}
  \end{prooftree}
}

\ex{$\vdash A \wedge (B \wedge C) \to B \wedge (C \wedge A)$ (lavagna)}{
  Ipotesi disponibile (da $\to I$): $[A \wedge (B \wedge C)]$.
  \begin{prooftree}
    \AxiomC{$[A \wedge (B \wedge C)]^1$}
    \RightLabel{$\wedge E_2$}
    \UnaryInfC{$B \wedge C$}
    \RightLabel{$\wedge E_1$}
    \UnaryInfC{$B$}
    \AxiomC{$[A \wedge (B \wedge C)]^1$}
    \RightLabel{$\wedge E_2$}
    \UnaryInfC{$B \wedge C$}
    \RightLabel{$\wedge E_2$}
    \UnaryInfC{$C$}
    \AxiomC{$[A \wedge (B \wedge C)]^1$}
    \RightLabel{$\wedge E_1$}
    \UnaryInfC{$A$}
    \RightLabel{$\wedge I$}
    \BinaryInfC{$C \wedge A$}
    \RightLabel{$\wedge I$}
    \BinaryInfC{$B \wedge (C \wedge A)$}
    \RightLabel{$\to I_1$}
    \UnaryInfC{$A \wedge (B \wedge C) \to B \wedge (C \wedge A)$}
  \end{prooftree}
  La stessa ipotesi si usa tre volte: un'ipotesi si può usare quante volte si
  vuole, e un solo $\to I$ scarica tutte le occorrenze.
}

\ex{II: $\vdash \varphi \to ((\varphi \to \bot) \to \bot)$}{
  \begin{prooftree}
    \AxiomC{$[\varphi]^2$}
    \AxiomC{$[\varphi \to \bot]^1$}
    \RightLabel{$\to E$}
    \BinaryInfC{$\bot$}
    \RightLabel{$\to I_1$}
    \UnaryInfC{$(\varphi \to \bot) \to \bot$}
    \RightLabel{$\to I_2$}
    \UnaryInfC{$\varphi \to ((\varphi \to \bot) \to \bot)$}
  \end{prooftree}
  Con l'abbreviazione $\neg\varphi$ per $\varphi \to \bot$ è la derivazione
  II$'$ di $\varphi \to \neg\neg\varphi$:
  \begin{prooftree}
    \AxiomC{$[\varphi]^2$}
    \AxiomC{$[\neg\varphi]^1$}
    \RightLabel{$\to E$}
    \BinaryInfC{$\bot$}
    \RightLabel{$\to I_1$}
    \UnaryInfC{$\neg\neg\varphi$}
    \RightLabel{$\to I_2$}
    \UnaryInfC{$\varphi \to \neg\neg\varphi$}
  \end{prooftree}
}

\ex{III: $\vdash (\varphi \to (\psi \to \sigma)) \to (\varphi \wedge \psi \to
\sigma)$}{
  \begin{prooftree}
    \AxiomC{$[\varphi \wedge \psi]^1$}
    \RightLabel{$\wedge E$}
    \UnaryInfC{$\psi$}
    \AxiomC{$[\varphi \wedge \psi]^1$}
    \RightLabel{$\wedge E$}
    \UnaryInfC{$\varphi$}
    \AxiomC{$[\varphi \to (\psi \to \sigma)]^2$}
    \RightLabel{$\to E$}
    \BinaryInfC{$\psi \to \sigma$}
    \RightLabel{$\to E$}
    \BinaryInfC{$\sigma$}
    \RightLabel{$\to I_1$}
    \UnaryInfC{$\varphi \wedge \psi \to \sigma$}
    \RightLabel{$\to I_2$}
    \UnaryInfC{$(\varphi \to (\psi \to \sigma)) \to (\varphi \wedge \psi \to
      \sigma)$}
  \end{prooftree}
  A lavagna è la stessa con $A, B, C$:
  $\vdash (A \to B \to C) \to A \wedge B \to C$, dove $A \to B \to C$ sta per
  $A \to (B \to C)$ (associatività a destra).
}

\ex{$A \to B,\ C \to \bot \vdash A \to (D \to B) \wedge (C \to D)$
(lavagna)}{
  Ipotesi: le due premesse (non scaricate), più $[A]$, $[D]$, $[C]$ dalle
  introduzioni di $\to$.
  \begin{prooftree}
    \AxiomC{$A \to B$}
    \AxiomC{$[A]^3$}
    \RightLabel{$\to E$}
    \BinaryInfC{$B$}
    \RightLabel{$\to I_1$}
    \UnaryInfC{$D \to B$}
    \AxiomC{$C \to \bot$}
    \AxiomC{$[C]^2$}
    \RightLabel{$\to E$}
    \BinaryInfC{$\bot$}
    \RightLabel{$\bot_E$}
    \UnaryInfC{$D$}
    \RightLabel{$\to I_2$}
    \UnaryInfC{$C \to D$}
    \RightLabel{$\wedge I$}
    \BinaryInfC{$(D \to B) \wedge (C \to D)$}
    \RightLabel{$\to I_3$}
    \UnaryInfC{$A \to (D \to B) \wedge (C \to D)$}
  \end{prooftree}
  Due cose da notare: $\to I_1$ scarica l'ipotesi $D$ che \textbf{non è mai
  stata usata} (scarico ``a vuoto'', lecito); nel ramo destro $D$ si ottiene
  con \textbf{ex falso}.
}

\ex{$\vdash \neg\neg A \to A$: serve RAA (lavagna)}{
  \textbf{Senza RAA} il tentativo fallisce ($\times$): dopo $\to I$ ho
  l'ipotesi $[(A \to \bot) \to \bot]$ e devo derivare $A$, ma per usare
  l'ipotesi con $\to E$ servirebbe $A \to \bot$, che non ho. In effetti
  $\neg\neg\varphi \to \varphi$ \textbf{non è derivabile nel sistema senza
  RAA} (Van Dalen).

  \textbf{Con RAA}: assumo anche $[A \to \bot]$, che RAA poi scarica.
  \begin{prooftree}
    \AxiomC{$[(A \to \bot) \to \bot]^2$}
    \AxiomC{$[A \to \bot]^1$}
    \RightLabel{$\to E$}
    \BinaryInfC{$\bot$}
    \RightLabel{RAA$_1$}
    \UnaryInfC{$A$}
    \RightLabel{$\to I_2$}
    \UnaryInfC{$\neg\neg A \to A$}
  \end{prooftree}
}

\ex{$((A \wedge B) \to \bot) \to \bot \vdash A$ (lavagna)}{
  Cioè $\neg\neg(A \wedge B) \vdash A$. La conclusione è un atomo, nessuna
  introduzione si applica: si parte con RAA, assumendo $[A \to \bot]$.
  \begin{prooftree}
    \AxiomC{$((A \wedge B) \to \bot) \to \bot$}
    \AxiomC{$[A \to \bot]^2$}
    \AxiomC{$[A \wedge B]^1$}
    \RightLabel{$\wedge E_1$}
    \UnaryInfC{$A$}
    \RightLabel{$\to E$}
    \BinaryInfC{$\bot$}
    \RightLabel{$\to I_1$}
    \UnaryInfC{$(A \wedge B) \to \bot$}
    \RightLabel{$\to E$}
    \BinaryInfC{$\bot$}
    \RightLabel{RAA$_2$}
    \UnaryInfC{$A$}
  \end{prooftree}
}

\ex{Esame 06/07/2021, es.\ 1: $\vdash (F \to G) \to \neg(F \wedge \neg G)$}{
  Scrivendo le negazioni con $\bot$: $(F \to G) \to ((F \wedge (G \to \bot))
  \to \bot)$. Due $\to I$ danno le ipotesi $[F \to G]$ e
  $[F \wedge (G \to \bot)]$; resta da derivare $\bot$.
  \begin{prooftree}
    \AxiomC{$[F \wedge (G \to \bot)]^1$}
    \RightLabel{$\wedge E_2$}
    \UnaryInfC{$G \to \bot$}
    \AxiomC{$[F \to G]^2$}
    \AxiomC{$[F \wedge (G \to \bot)]^1$}
    \RightLabel{$\wedge E_1$}
    \UnaryInfC{$F$}
    \RightLabel{$\to E$}
    \BinaryInfC{$G$}
    \RightLabel{$\to E$}
    \BinaryInfC{$\bot$}
    \RightLabel{$\to I_1$}
    \UnaryInfC{$F \wedge \neg G \to \bot$}
    \RightLabel{$\to I_2$}
    \UnaryInfC{$(F \to G) \to \neg(F \wedge \neg G)$}
  \end{prooftree}
  \textbf{Controllo semantico (lavagna)}: cerchiamo un controesempio. Perché
  la formula sia falsa serve $F \to G = 1$ e $\neg(F \wedge \neg G) = 0$, cioè
  $F = 1$ e $G = 0$; ma allora $F \to G = 0$. \textbf{Nessun controesempio}:
  la formula è una tautologia, quindi (per completezza) una dimostrazione
  esiste.
}

\ex{Una formula non dimostrabile: $(F \to G) \to (\neg F \to \neg G)$}{
  Il tentativo bottom-up porta a tre ipotesi, $[F \to G]$, $[F \to \bot]$,
  $[G]$, da cui bisognerebbe derivare $\bot$: \textbf{nessuna dimostrazione}
  ($\times$). Per esserne sicuri si cerca un \textbf{controesempio}:
  \[
    F = 0,\ G = 1: \qquad \underbrace{F \to G}_{1} \to
    \underbrace{(\underbrace{\neg F}_{1} \to \underbrace{\neg G}_{0})}_{0}
    \;=\; 0.
  \]
  La formula non è una tautologia, quindi (per correttezza) \textbf{non è
  derivabile}: $\not\vdash (F \to G) \to (\neg F \to \neg G)$. È l'errore
  ``negazione dell'antecedente'' (la contrapposizione corretta è
  $\neg G \to \neg F$).
}

\nt{
  \textbf{Metodo d'esame.} Per ``dimostrare o confutare'' $\Gamma \vdash
  \varphi$:
  \begin{itemize}
    \item per mostrare che \emph{vale} si esibisce un albero di derivazione;
    \item per mostrare che \emph{non vale} si esibisce un'interpretazione che
          rende vere le formule di $\Gamma$ e falsa $\varphi$ (un
          controesempio): per correttezza, se $\Gamma \not\models \varphi$
          allora $\Gamma \not\vdash \varphi$.
  \end{itemize}
  Conviene cercare prima il controesempio: se non c'è, si sa che la
  derivazione esiste.
}

\subsection{Le derivazioni sono alberi}

Ogni derivazione è un \textbf{albero}: le \textbf{foglie} contengono le
premesse (ipotesi), la \textbf{radice} è la conclusione. Formalmente:

\dfn{Derivazioni (Van Dalen 1.4.1)}{
  L'insieme delle derivazioni è il più piccolo insieme $X$ tale che:
  \begin{enumerate}
    \item[(1)] l'albero con un solo elemento $\varphi$ appartiene a $X$, per
          ogni $\varphi \in \PROP$;
    \item[(2$\wedge$)] se
          $\begin{array}{c}\mathcal{D}\\ \varphi\end{array}$,
          $\begin{array}{c}\mathcal{D}'\\ \varphi'\end{array} \in X$, allora
          \AxiomC{$\mathcal{D}$}
          \noLine
          \UnaryInfC{$\varphi$}
          \AxiomC{$\mathcal{D}'$}
          \noLine
          \UnaryInfC{$\varphi'$}
          \BinaryInfC{$\varphi \wedge \varphi'$}
          \DisplayProof
          $\in X$; \quad se
          $\begin{array}{c}\mathcal{D}\\ \varphi \wedge \psi\end{array} \in X$,
          allora
          \AxiomC{$\mathcal{D}$}
          \noLine
          \UnaryInfC{$\varphi \wedge \psi$}
          \UnaryInfC{$\varphi$}
          \DisplayProof
          ,
          \AxiomC{$\mathcal{D}$}
          \noLine
          \UnaryInfC{$\varphi \wedge \psi$}
          \UnaryInfC{$\psi$}
          \DisplayProof
          $\in X$;
    \item[(2$\to$)] se
          $\begin{array}{c}\varphi\\ \mathcal{D}\\ \psi\end{array} \in X$,
          allora
          \AxiomC{$[\varphi]$}
          \noLine
          \UnaryInfC{$\mathcal{D}$}
          \noLine
          \UnaryInfC{$\psi$}
          \UnaryInfC{$\varphi \to \psi$}
          \DisplayProof
          $\in X$; \quad se
          $\begin{array}{c}\mathcal{D}\\ \varphi\end{array}$,
          $\begin{array}{c}\mathcal{D}'\\ \varphi \to \psi\end{array} \in X$,
          allora
          \AxiomC{$\mathcal{D}$}
          \noLine
          \UnaryInfC{$\varphi$}
          \AxiomC{$\mathcal{D}'$}
          \noLine
          \UnaryInfC{$\varphi \to \psi$}
          \BinaryInfC{$\psi$}
          \DisplayProof
          $\in X$;
    \item[(2$\bot$)] se
          $\begin{array}{c}\mathcal{D}\\ \bot\end{array} \in X$, allora
          \AxiomC{$\mathcal{D}$}
          \noLine
          \UnaryInfC{$\bot$}
          \UnaryInfC{$\varphi$}
          \DisplayProof
          $\in X$; \quad se
          $\begin{array}{c}\neg\varphi\\ \mathcal{D}\\ \bot\end{array} \in X$,
          allora
          \AxiomC{$[\neg\varphi]$}
          \noLine
          \UnaryInfC{$\mathcal{D}$}
          \noLine
          \UnaryInfC{$\bot$}
          \UnaryInfC{$\varphi$}
          \DisplayProof
          $\in X$.
  \end{enumerate}
}

\nt{
  È di nuovo una definizione induttiva (come $\PROP$): quindi si possono
  dimostrare proprietà \textbf{per induzione sulle derivazioni} e definire
  funzioni per ricorsione su di esse. È lo strumento con cui si dimostra la
  correttezza.
}

\dfn{Derivabilità e teoremi (Van Dalen 1.4.2)}{
  Scriviamo
  \[ \Gamma \vdash \varphi \]
  se esiste una derivazione con conclusione $\varphi$ e con tutte le ipotesi
  \textbf{non scaricate} in $\Gamma$. Si dice che $\varphi$ è
  \textbf{derivabile} da $\Gamma$; il simbolo $\vdash$ si chiama
  \emph{turnstile}.
  Se $\Gamma = \emptyset$ scriviamo $\vdash \varphi$ e diciamo che $\varphi$
  è un \textbf{teorema}.
}

\nt{
  $\Gamma$ può contenere ipotesi superflue: se $\Gamma \vdash \varphi$ e
  $\Gamma \subseteq \Gamma'$ allora $\Gamma' \vdash \varphi$. Gli esempi I,
  II, III sopra hanno tutte le ipotesi scaricate: le loro conclusioni sono
  teoremi.
}

% =====================================================================
\section{Correttezza e completezza}
% =====================================================================

\thm{Teorema di completezza (in senso lato)}{
  Le nozioni di derivabilità e di verità \textbf{coincidono}:
  \[ \Gamma \vdash \varphi \iff \Gamma \models \varphi. \]
}

\cor{}{
  L'insieme dei \textbf{teoremi} coincide con l'insieme delle
  \textbf{tautologie}: $\vdash \varphi \iff \models \varphi$.
}

Le due direzioni hanno nomi diversi.

\thm{Correttezza (Soundness, Van Dalen 1.5.1)}{
  \[ \Gamma \vdash \varphi \ \Rightarrow\ \Gamma \models \varphi. \]
}

\pf{Idea della dimostrazione}{
  Per induzione sulle derivazioni: si mostra che per ogni derivazione
  $\mathcal{D}$ con conclusione $\varphi$ e ipotesi in $\Gamma$ vale
  $\Gamma \models \varphi$. Caso base: l'albero di un solo elemento $\varphi$
  ha $\varphi \in \Gamma$, quindi $\Gamma \models \varphi$. Passo: ogni regola
  conserva la verità. Ad esempio per $\to I$: per ipotesi induttiva
  $\Gamma, \varphi \models \psi$, quindi per il teorema di deduzione semantico
  $\Gamma \models \varphi \to \psi$; per RAA: da $\Gamma, \neg\varphi \models
  \bot$ segue che nessuna valutazione rende vero $\Gamma$ e falsa $\varphi$,
  cioè $\Gamma \models \varphi$. (Dettagli sul libro.)
}

\thm{Completezza (Van Dalen 1.5.13)}{
  \[ \Gamma \models \varphi \ \Rightarrow\ \Gamma \vdash \varphi. \]
}

\pf{Idea della dimostrazione}{
  Si dimostra la contronominale $\Gamma \not\vdash \varphi \Rightarrow \Gamma
  \not\models \varphi$, mostrando che:
  \begin{center}
    $\Gamma \not\vdash \varphi$ \quad sse \quad $\Gamma \cup \{\neg\varphi\}$
    è consistente \quad sse \quad esiste una valutazione $v$ tale che
    $\sem{\psi}_v = 1$ per ogni $\psi \in \Gamma \cup \{\neg\varphi\}$.
  \end{center}
  Una tale $v$ rende vero $\Gamma$ e falsa $\varphi$: è un controesempio,
  quindi $\Gamma \not\models \varphi$. La prima equivalenza usa il lemma
  seguente; la seconda (``ogni insieme consistente ha un modello'') è la parte
  difficile, e si ottiene estendendo l'insieme a uno \emph{massimamente
  consistente} (vedi libro).
}

\dfn{Consistenza (sintattica)}{
  \begin{itemize}
    \item $\Gamma$ è \textbf{consistente} se $\Gamma \not\vdash \bot$;
    \item $\Gamma$ è \textbf{inconsistente} se $\Gamma \vdash \bot$.
  \end{itemize}
}

\mlenma{Un lemma utile (Van Dalen 1.5.5)}{
  \begin{enumerate}[label=(\alph*)]
    \item $\Gamma \cup \{\neg\varphi\}$ è inconsistente
          $\ \Rightarrow\ \Gamma \vdash \varphi$;
    \item $\Gamma \cup \{\varphi\}$ è inconsistente
          $\ \Rightarrow\ \Gamma \vdash \neg\varphi$.
  \end{enumerate}
  Detto altrimenti (contronominale di (a)): se $\Gamma \not\vdash \varphi$
  allora $\Gamma \cup \{\neg\varphi\}$ è consistente.
}

\pf{Dimostrazione}{
  Le ipotesi di (a) e (b) danno due derivazioni di $\bot$, con ipotesi in
  $\Gamma \cup \{\neg\varphi\}$ e in $\Gamma \cup \{\varphi\}$
  rispettivamente. Applicando RAA alla prima e $\to I$ alla seconda si
  scarica l'ipotesi in più:
  \begin{center}
    \AxiomC{$[\neg\varphi]$}
    \noLine
    \UnaryInfC{$\mathcal{D}$}
    \noLine
    \UnaryInfC{$\bot$}
    \RightLabel{RAA}
    \UnaryInfC{$\varphi$}
    \DisplayProof
    \qquad\qquad
    \AxiomC{$[\varphi]$}
    \noLine
    \UnaryInfC{$\mathcal{D}'$}
    \noLine
    \UnaryInfC{$\bot$}
    \RightLabel{$\to I$}
    \UnaryInfC{$\neg\varphi$}
    \DisplayProof
  \end{center}
  Le derivazioni ottenute hanno conclusione $\varphi$ (risp.\ $\neg\varphi$)
  e ipotesi in $\Gamma$.
}

\nt{
  \textbf{Che cosa ci dice la completezza.}
  \begin{itemize}
    \item Il compito ``creativo'' di costruire derivazioni può essere
          sostituito da quello (noioso ma automatico) di controllare
          tautologie con le tavole di verità, e viceversa.
    \item La consistenza \emph{sintattica} ($\Gamma \not\vdash \bot$) coincide
          con la soddisfacibilità \emph{semantica} (esiste un modello): è il
          legame ``ha un modello $\iff$ consistente'' anticipato nel capitolo
          sulla verità.
    \item Negli esercizi: per mostrare $\Gamma \vdash \varphi$ si dà una
          derivazione; per mostrare $\Gamma \not\vdash \varphi$ si dà un
          controesempio semantico (correttezza).
  \end{itemize}
}
